% ch-governing.tex -- theory: the long-wave equations, their scaling and their limits
\chapter{The governing equations}
\label{ch:governing}

This chapter states what Yuzen actually solves: one evolution equation for the
film thickness, one transport equation for the solute, the approximation that
reduces the Navier--Stokes equations to them, and the scaling that turns them
into the dimensionless form the solver works in. Read it if you want to know
what a given input key does to the mathematics, or where the model stops being
trustworthy. The individual flux terms are taken one at a time in
\cref{ch:terms}; the keys themselves are in \cref{ch:ref-physics,ch:ref-domain}.

\section{The long-wave approximation}
\label{sec:governing:longwave}

When the film thickness $h$ is small compared with the representative in-plane
length $L$, the Navier--Stokes equations collapse into a single evolution
equation for $h(x,y,t)$. This is the long-wave --- or lubrication ---
approximation, and it is the equation Yuzen solves. Four assumptions go into
it, and each of them buys something and costs something.

\begin{itemize}
  \item \textbf{The aspect ratio is small}, $\varepsilon = h_0/L_0 \ll 1$. This
    is what allows the cross-film coordinate to be eliminated altogether: you
    solve a problem in one or two dimensions instead of two or three, and the
    unknown is a surface rather than a volume. The cost is that every
    conclusion is an expansion in $\varepsilon$, so nothing the model says
    about a feature as short as the film is thick can be believed.
  \item \textbf{Inertia is negligible.} The velocity profile across the film is
    then exactly parabolic, with no slip at the substrate (or a Navier slip
    length, \key{Physics.Slip}: \cref{sec:terms:slip}) and a shear stress at
    the free surface set solely by the surface-tension gradient. This is what
    makes the depth-integrated flux $\Qvec$ an explicit algebraic function of
    $h$ and the driving pressures, with no momentum equation left to solve. The
    cost is that fast rotation and low-viscosity films, where the profile
    departs from the parabola, are outside the model.
  \item \textbf{The pressure is uniform across the film.} Capillary pressure,
    gravity and disjoining pressure all enter one single pressure. This is what
    lets every physical effect be added as one more term in the same bracket,
    and it is why switching a term off is a matter of zeroing a coefficient
    rather than changing the equation. The cost is that no effect with
    structure through the thickness --- a pressure that varies with depth ---
    can be represented at all.
  \item \textbf{The interface slope is small}, so the curvature may be
    linearized to $\nabla^2 h$; that is, the denominator of the exact mean
    curvature $\nabla\cdot\bigl(\nabla h/\sqrt{1+|\nabla h|^2}\bigr)$ is set to
    1. This is what keeps the capillary term linear in the highest derivative
    and the whole equation tractable. The cost is a direct error wherever the
    surface is steep, which is exactly where a contact line is; see
    \cref{sec:governing:breakdown}.
\end{itemize}

\section{The dimensional equations}
\label{sec:governing:dimensional}

In dimensional variables the film thickness obeys

\begin{equation}\label{eq:height-dim}
  \frac{\partial h}{\partial t} = -\nabla\cdot\Qvec \;-\; E(h,c) \;+\; S(x,t),
\end{equation}

where $\Qvec$ is the flux integrated across the thickness, $E$ is the rate at
which evaporation removes thickness, and $S$ is liquid supplied from a nozzle.
The flux collects six contributions:

\begin{equation}\label{eq:flux-dim}
  \begin{aligned}
    \Qvec = {}& \frac{\sigma h^{3}}{3\mu}\,\nabla\nabla^{2}(h+s)
      && \text{capillary pressure} \\
      &{}- \frac{\rho g_z h^{3}}{3\mu}\,\nabla(h+s)
      && \text{vertical gravity} \\
      &{}+ \frac{h^{3}}{3\mu}\,\nabla\Pi(h)
      && \text{disjoining pressure} \\
      &{}+ \frac{h^{2}}{2\mu}\,\nabla\sigma
      && \text{Marangoni stress} \\
      &{}+ \frac{\rho g_\parallel h^{3}}{3\mu}
      && \text{in-plane gravity} \\
      &{}+ \frac{\rho \omega^{2} h^{3}}{3\mu}\,r\,\er
      && \text{centrifugal force.}
  \end{aligned}
\end{equation}

Here $h$ is the film thickness and $s(x,y)$ is the height of the substrate, so
the free surface sits at $h+s$. The distinction matters and it is not
cosmetic. \textbf{Capillary pressure and vertical gravity are driven by the
shape of the free surface $h+s$}, because both are consequences of where the
liquid surface is in space; \textbf{disjoining pressure is a function of the
film thickness $h$ alone}, because it describes an interaction between the
substrate and the free surface whose strength depends only on how far apart
they are. One practical consequence is that a substrate shape does nothing
unless at least one of the first two terms is enabled
(\cref{sec:terms:vgravity}).

Both $\mu = \mu(c)$ and $\sigma = \sigma(c)$ depend on the solute
concentration (\cref{sec:conc:properties}), which is what couples
\cref{eq:height-dim} to the transport equation below.

\begin{note}[The velocity Yuzen writes out is not a new unknown]
  Nothing in \cref{eq:height-dim,eq:flux-dim} solves for a velocity. The
  in-plane velocity in the field output is the quantity $\Qvec/h$ --- the
  depth-averaged velocity implied by the flux, summed over the active terms.
  The \file{.vts} files of the well-mixed models also report it term by term,
  one array per contribution in \cref{eq:flux-dim} except the disjoining
  pressure (\cref{sec:output:vtk}). That decomposition is the most direct way
  to see which physical effect is actually moving the liquid at a given place
  and time.
\end{note}

\section{Solute transport}
\label{sec:governing:concentration}

The solute (resin) concentration $c$ is a mass fraction, so $c = 1$ means
completely dry. In the well-mixed model it is taken to be uniform across the
thickness, and what is conserved is the solute held in a column of the film,
$hc$:

\begin{equation}\label{eq:conc-dim}
  \frac{\partial (hc)}{\partial t} + \nabla\cdot\bigl(c\,\Qvec\bigr)
    = \nabla\cdot\bigl(D(c)\,h\,\nabla c\bigr) \;+\; c_S\,S .
\end{equation}

The solute is carried by the same flux $\Qvec$ that moves the liquid in
\cref{eq:height-dim}, spreads by diffusion along the film, and arrives with the
liquid a nozzle supplies. $c_S$ is the concentration that liquid arrives at,
chosen by \key{Physics.Source.ConcType} (\cref{sec:terms:source}): $c_S = c$,
the concentration already present where it lands (\val{"local"}, the
default), or $c_S = \key{Physics.Source.Conc}$ (\val{"supplied"}), mixed
through the thickness of the column it lands on. \textbf{Evaporation does not appear}: it removes
solvent and leaves the solute behind. This is the equation Yuzen evolves ---
the unknown it advances is the solute $hc$ of each cell, with the film $h$ of
the same step --- so the solute of a closed film is conserved to rounding and a
uniform $c$ stays uniform whatever the film does, unless liquid of another
concentration is supplied (\cref{sec:disc:wm-conservative}).

Subtracting $c$ times \cref{eq:height-dim} gives the same law for $c$ itself,
with $\uvec = \Qvec/h$ the mixed velocity of all the enabled terms:
%
\begin{equation}\label{eq:conc-adv}
  \frac{\partial c}{\partial t} + \uvec\cdot\nabla c
    = \frac{1}{h}\,\nabla\cdot\bigl(D(c)\,h\,\nabla c\bigr)
      \;+\; \frac{E c}{h} \;+\; \frac{S\,(c_S - c)}{h}.
\end{equation}

It is the clearer form to read, and its evaporation term is the one that makes drying
interesting. Evaporation removes
solvent but not solute, so whatever is left behind becomes more concentrated;
the rate at which it does so is the thickness lost per unit time, $E$, divided
by the thickness that is there to dilute it, $h$. \textbf{Because it is divided
by $h$, concentration rises fastest where the film is thinnest.} That is the
mechanism behind edge effects in a drying coating, and it is also why a
contact-line calculation needs the precursor-film correction of
\cref{sec:terms:evaporation}: without it, the thinnest region in the domain
would dry out first and fastest. The last term is the nozzle's, and it is 0
when the liquid arrives at the local concentration; liquid supplied at
another $c_S$ pulls $c$ toward $c_S$ at the rate $S/h$, so a thin film under
the nozzle takes on the supplied concentration quickly and a thick one slowly.

If you need the distribution of concentration through the thickness rather than
a depth average, use the \val{"CZ"} model instead (\cref{sec:conc:cz}).

\section{Where the assumptions break down}
\label{sec:governing:breakdown}

\paragraph{Large slope.} Two separate errors appear together. The linearized
curvature loses accuracy, and so does the small-angle substitution
$1-\cos\thetae \approx \thetae^{2}/2$ used to calibrate the disjoining pressure
(\cref{sec:wetting:dp}). The second error is easy to quantify exactly:

\begin{yzntable}{Error of $1-\cos\thetae \approx \thetae^{2}/2$\label{tab:cos-error}}{@{}d d@{}}
\textsf{\bfseries $\thetae$ [deg]} & \textsf{\bfseries error} \\ \midrule
10 & 0.3 \% \\
20 & 1.0 \% \\
30 & 2.3 \% \\
45 & 5.3 \% \\
90 & 23 \% \\
\end{yzntable}

\key{Physics.DisjoiningPressure.ContactAngle} accepts anything in $[0:90]$
degrees, but if you intend to use the angle quantitatively, keep it small.

\paragraph{Inertia.} At high rotation rates, or in a low-viscosity film, the
velocity profile departs from the parabola the derivation assumes. Nothing in
the solver detects this; the equations remain perfectly well behaved and simply
describe a different problem from the one you have.

\paragraph{In-plane lengths approaching the film thickness.} The
neighborhood of a contact line is by construction a region of large slope and
short in-plane length --- the two things the approximation is least able to
handle. The disjoining pressure is the device that lets the long-wave equation
be carried into that region at all, and \cref{ch:wetting} is about what it can
and cannot deliver there.

All of this is concentrated in one number. The assumption that the film is thin
appears in the dimensionless group

\[
  \As = \left(\frac{L_0}{h_0}\right)^{2},
\]

and a configuration in which $\As$ falls toward order 1 is outside the range
of the approximation altogether.

\section{Scales}
\label{sec:governing:scales}

Every value in the input file is dimensional --- meters, seconds, pascal
seconds --- but the solver works entirely in dimensionless variables. Five
scales do the conversion.

\begin{yzntable}{The scales used for non-dimensionalization\label{tab:scales}}{@{}l L{0.42\textwidth} L{0.44\textwidth}@{}}
\textsf{\bfseries Symbol} & \textsf{\bfseries What it is} &
\textsf{\bfseries Set by} \\ \midrule
$L_0$ & In-plane length: the capillary length $\sqrt{\sigma_0/(\rho g)}$, with $g = \SI{9.80665}{\metre\per\second\squared}$ & \key{LengthScaleType}; with \val{"custom"}, $L_0$ is taken from \key{LengthScale} instead \\
$h_0$ & Film thickness & \key{Initial.SurfaceShape.TotalHeight} \\
$T_0$ & Time, $\mu_0 L_0^{4}/(\sigma_0 h_0^{3})$ & follows from $\mu_0$, $\sigma_0$, $L_0$ and $h_0$ \\
$U_0$ & Velocity, $L_0/T_0$ & follows from $L_0$ and $T_0$ \\
$\varepsilon$ & Aspect ratio, $h_0/L_0$ & follows from $h_0$ and $L_0$ \\
\end{yzntable}

The reference properties are $\sigma_0 =$ \key{Fluid.SurfaceTension.BaseValue},
$\rho =$ \key{Fluid.Density.Value} and $\mu_0 =$
\key{Fluid.Viscosity.BaseValue}. Writing \key{LengthScaleType = "custom"}
replaces the capillary length by whatever \key{LengthScale} says.

The dimensionless thickness is $H = h/h_0$ and the dimensionless in-plane
coordinate is $X = x/L_0$. It follows that

\begin{equation}\label{eq:slope-conversion}
  \text{physical slope} \;=\; \varepsilon \times \text{dimensionless slope},
\end{equation}

and \cref{eq:slope-conversion} is not a footnote: \textbf{the entire discussion
of contact angles in \cref{ch:wetting} rests on it}. A dimensionless profile
that rises by 1 over a unit of $X$ is a physical surface inclined at
$\varepsilon$, not at $45^\circ$. Four of the five --- $L_0$, $h_0$, $T_0$ and
$U_0$ --- are printed in the run log on standard output under the heading
\code{Scaling factors} (\cref{sec:output:running}); read them rather than
recomputing them. $\varepsilon$ is not printed, but it is $h_0/L_0$ of the two
that are.

\section{The dimensionless groups}
\label{sec:governing:groups}

Twelve dimensionless groups carry the whole of the input into the equations.

\begin{yzntable}{The dimensionless groups\label{tab:groups}}{@{}l L{0.225\textwidth} L{0.22\textwidth} L{0.375\textwidth}@{}}
\textsf{\bfseries Symbol} & \textsf{\bfseries Definition} &
\textsf{\bfseries Meaning} & \textsf{\bfseries Set by} \\ \midrule
$\As$   & $(L_0/h_0)^{2}$ & Inverse aspect ratio squared; the strength of capillary pressure & \key{Initial.SurfaceShape.TotalHeight} together with $L_0$ \\
$\Bo$   & $\rho g_z L_0^{2}/\sigma_0$ & Bond number; gravity normal to the substrate & \key{Physics.\bk VerticalGravity.\bk Acceleration} \\
$\Pgx$  & $\rho g_x L_0^{2}/(\sigma_0\varepsilon)$ & In-plane gravity along $x$; non-zero on a tilted substrate & \key{Physics.\bk HorizontalGravity.\bk Acceleration[0]} \\
$\Pgy$  & $\rho g_y L_0^{2}/(\sigma_0\varepsilon)$ & In-plane gravity along $y$ & \key{Physics.\bk HorizontalGravity.\bk Acceleration[1]} \\
$\Spn$  & $\rho \omega^{2} L_0^{4}/(\sigma_0 h_0)$ & Spin number; centrifugal force & \key{Physics.Centrifugal.RotationRate} \\
$\Ma$   & $(\mathrm{d}\sigma/\mathrm{d}c)/\sigma_0$ & Marangoni number; how strongly surface tension depends on concentration & \key{Fluid.SurfaceTension.Coef} \\
$\Dpn$  & $\displaystyle\frac{\As\,(n-1)(m-1)\,\thetae^{2}}{2\,\Hstar\,(n-m)}$ & Disjoining-pressure coefficient, fixed by the target equilibrium contact angle & \key{Physics.DisjoiningPressure.*} together with \key{PrecursorRatio} \\
$\Hstar$ & \key{PrecursorRatio} & Dimensionless precursor thickness; the root of $\Pihat(\Hstar) = 0$ & \key{Initial.SurfaceShape.PrecursorRatio} \\
$\Ev$   & $E_0 T_0/h_0$ & Evaporation number; how fast evaporation thins the film & \key{Physics.Evaporation.Coef} \\
$\Pen$  & $E_0 h_0/D_0$ & P\'eclet number; evaporation against cross-film diffusion & \key{Physics.Evaporation.Coef}, \key{Fluid.Diffusivity.BaseValue} \\
$\Df$   & $D_0 T_0/L_0^{2}$ & Diffusion number; the speed of in-plane diffusion & \key{Fluid.Diffusivity.BaseValue} \\
$\Mr$   & $M_{\mathrm{resin}}/M_{\mathrm{solvent}}$ & Molar-mass ratio, used by Raoult's law & \key{Fluid.MolarMass.Solvent}, \key{Fluid.MolarMass.Resin} \\
$\Lambda$ & $\lambda/h_0$ & Dimensionless slip length; mobility $H^3/3 + \Lambda H^2$ (\cref{sec:terms:slip}) & \key{Physics.Slip.Length} \\
\end{yzntable}

In these definitions $g_z$, $g_x$ and $g_y$ are the components of gravitational
acceleration given in the input in \si{\metre\per\second\squared}, $\omega$ is
the rotation rate in \si{\radian\per\second}, $E_0$ is the evaporation
coefficient in \si{\metre\per\second} and $D_0$ is the reference diffusivity in
\si{\metre\squared\per\second}.

$\Bo$, $\Pgx$, $\Pgy$, $\Spn$, $\Ma$ and $\Dpn$ are \textbf{set to zero} when
the corresponding term is disabled in the input, and the term is then not
assembled at all (\cref{sec:repro:order}). $\Ev$, $\Pen$ and $\Df$ always carry a value, and evaporation and
diffusion are switched on and off by separate keys instead. Note the
side-effect of the first rule: $\Ma = 0$ removes the Marangoni stress and, at
the same time, removes the concentration dependence of surface tension entirely
(\cref{sec:conc:properties}).

One number is worth remembering. If you leave the length scale at the capillary
length and set \key{Acceleration = 9.80665} --- that is,
$g = \SI{9.80665}{\metre\per\second\squared}$ --- then $\Bo = 1$ exactly, since
the capillary length is defined as the length at which gravity and surface
tension balance. \key{Acceleration} is therefore the key to change when you want to
move $\Bo$.

\section{The dimensionless equations}
\label{sec:governing:nondim}

With those scales and groups, \cref{eq:height-dim,eq:flux-dim} become

\begin{equation}\label{eq:height}
  \frac{\partial H}{\partial T} = -\nabla\cdot\Qvec - \Ev\,\hat{E}(H,C) + \hat{S},
\end{equation}

\begin{equation}\label{eq:flux}
  \Qvec = \frac{H^{3}}{3\hat{\mu}}
    \Bigl[\,\hat{\sigma}\,\nabla\nabla^{2}(H+S)
      - \Bo\,\nabla(H+S)
      + \Dpn\,\nabla\Pihat(H)
      + \Pvec
      + \Spn\,r\,\er \Bigr]
    + \frac{\As\,\Ma}{2}\,\frac{H^{2}}{\hat{\mu}}\,\nabla C,
\end{equation}

and the solute balance \cref{eq:conc-dim} becomes

\begin{equation}\label{eq:conc}
  \frac{\partial (HC)}{\partial T} + \nabla\cdot\bigl(C\,\Qvec\bigr)
    = \Df\,\nabla\cdot\bigl(\hat{D}(C)\,H\,\nabla C\bigr) + C_S\,\hat{S},
\end{equation}

with $C_S = c_S$ the concentration the supplied liquid arrives at ($C$ itself
by default), and the disjoining pressure

\begin{equation}\label{eq:pihat}
  \Pihat(H) = (\Hstar/H)^{n} - (\Hstar/H)^{m}.
\end{equation}

Here $\hat{\mu}(C)$, $\hat{\sigma}(C)$ and $\hat{D}(C)$ are the viscosity,
surface tension and diffusivity normalized to 1 at the reference concentration, and $\Pvec = (\Pgx, \Pgy)$.
Note that $S$ without a hat is the dimensionless substrate height, while
$\hat{S}$ with one is the dimensionless source; they are different quantities
that unfortunately share a letter.

The single most useful thing to notice in \cref{eq:flux} is what makes it
vanish. Five of the six contributions sit inside one bracket multiplied by
$H^{3}/3\hat{\mu}$, so $\Qvec = 0$ requires the contents of that bracket to be
\emph{uniform in space}. That statement --- uniformity of the bracket, not
smallness of any individual term --- is the equilibrium condition, and it is
where \cref{ch:wetting} starts. In the discretization, thickness and
concentration live at cell centers and velocities on cell faces
(\cref{ch:discretization}).
